\clearpage
\section{Limitations}
\label{appendix: limitations}

\benchmarkname is the first comprehensive benchmark designed to evaluate the long-context document understanding capabilities of LVLMs. While our benchmark addresses significant gaps in the previous datasets, we acknowledge several limitations.

One primary limitation is the scale of the benchmark. Currently, our benchmark includes a test set comprising 135 documents and 1,082 questions. It is much smaller compared to previous datasets. The complexity and difficulty of annotations limit the scale of our benchmark. As a long-context benchmark, our documents average about 50 pages and 20,000 tokens. And most questions require either complicated reasoning or cross-page comprehension. It takes more than one hour for an expert-level annotator to read through a single document, and then edit existing instances and create new instances on this document. Given the purpose of \benchmarkname as an evaluation benchmark, we prioritize annotation quality over quantity. Moreover, the results presented in Sections 3.3 and 3.4 confirm that the scale of our benchmark is sufficient for fine-grained evaluations across different document types, evidence sources, evidence pages, \etc. Additionally, we plan to expand our benchmark by adding more documents and questions in future iterations.

We roughly categorize these questions into three types, \ie single-page, cross-page, and unanswerable questions, based on whether evidence can be found in the documents and the number of evidence pages. However, unlike MMBench~\cite{mmbench2023} or MathVista~\cite{lu2024mathvista}, we provide no further taxonomy to classify some (\eg 7 or 20) fine-grained, evaluated reasoning or perception capabilities out of two main reasons: (1) Prior (\ie pre-annotation) taxonomy limits the diversity of the questions. Therefore we provide no predefined classifications in our guideline and encourage the expert-level annotators to freely write questions without constraints. (2) The intrinsic complexity of document understanding presents significant challenges for establishing a posterior (\ie post-annotation) taxonomy.

While there exist limitations in our benchmark, \benchmarkname surely represents a significant step forward in this field. We would iteratively maintain and refine this benchmark and hope it could push forward the development of long-context document understanding.